\documentclass[12pt]{article}

\input{../../_assets/preambulo.tex}

\newcommand{\E}{\mathrm{E}}
\newcommand{\V}{\mathrm{Var}}
\newcommand{\R}{\mathbb{R}}
\newcommand{\Rmas}{\R^{+}}
\newcommand{\Ss}{\mathbb{S}}
\newcommand{\Z}{\mathbb{Z}}
\newcommand{\X}{\mathcal{X}}
\newcommand{\N}{\mathbb{N}}
\newcommand{\I}{\mathrm{Im}}
\newcommand{\Rmenos}{\R^{-}}
\newcommand{\norma}[1]{\lVert #1 \rVert}
\newcommand{\prodesc}[2]{\langle #1, #2 \rangle}
\newcommand{\normagenerica}{\lVert \cdot \rVert}
\newcommand{\Rceromas}{\R_0^{+}}
\newcommand{\tr}{\text{tr}}

\usepackage{pgfplots}
\usepgfplotslibrary{fillbetween}

\begin{document}

    % 1. Foto de fondo
    % 2. Título
    % 3. Encabezado Izquierdo
    % 4. Color de fondo
    % 5. Coord x del titulo
    % 6. Coord y del titulo
    % 7. Fecha

    
    \input{../../_assets/portada}
    \portadaExamen{ffccA4.jpg}{Curvas y \\ Superficies \\Examen IV}{Curvas y Superficies. Examen IV}{MidnightBlue}{-8}{28}{2026}{José Manuel Sánchez Varbas}

    \begin{description}
        \item[Asignatura] Curvas y Superficies.
        \item[Curso Académico] 2023-24.
        \item[Grado] Grado en Matemáticas.
        \item[Descripción] Ordinaria.
        \item[Fecha] 17 de Junio de 2024.  
    \end{description}
    \newpage


    % ------------------------------------

    \begin{ejercicio}~
        \begin{enumerate}
            \item[a)] Sea $\alpha : (0,1) \to \mathbb{R}^3$ una curva diferenciable p.p.a. y tal que su traza está contenida en la esfera centrada en el origen de radio $r > 0$. Comprueba que la curvatura de $\alpha$ es siempre mayor o igual que $1/r$.

            \item[b)] Si $\beta : (0,1) \to \mathbb{R}^3$ una curva diferenciable p.p.a. con curvatura no nula en todo instante, entonces $\gamma : (0,1) \to \mathbb{R}^3$, $\gamma(t)=\beta'(t)$, es una curva regular con curvatura mayor o igual que $1$.
        \end{enumerate}
    \end{ejercicio}

    \begin{ejercicio}
    Sea $S=\{X(u,v)\mid u\in (0,+\infty), v\in (-\frac{\pi}{2},\frac{\pi}{2})\}$ donde
    $$
    X(u,v)=(u\cos(v),u\sen(v),\ln(\cos(v))+u).
    $$

    Se pide:
    \begin{enumerate}
        \item[a)] Demuestra que $X$ es una parametrización global de $S$ y que $S$ es una superficie orientable.

        \item[b)] Calcula la primera y segunda formas fundamentales de $S$. Prueba que el punto $p=(1,0,1)\in S$ y determina el plano tangente de la superficie en dicho punto.

        \item[c)] Calcula la curvatura de Gauss y la curvatura media de $S$ y clasifica los puntos de $S$ en elípticos, hiperbólicos, llanos $($parabólicos y planos$)$.

        \item[d)] Prueba que $S$ es una superficie reglada. Razonar si $S$ es localmente isométrica o no a un plano.

        \item[e)] Sea $\pi:S\to \{z=0\}$ la aplicación dada por $\pi(x,y,z)=(x,y,0)$. Prueba que $\pi$ es un difeomorfismo sobre su imagen.
    \end{enumerate}
    \end{ejercicio}

    \begin{ejercicio}
    Razonar si son verdaderas o falsas las siguientes afirmaciones:

    \begin{enumerate}
        \item[a)] Si $S_1,S_2$ son superficies regulares orientables con aplicaciones de Gauss $N_1:S_1\to \mathbb{S}^2$ y $N_2:S_2\to \mathbb{S}^2$ y $f:S_1\to S_2$ es un difeomorfismo local tal que $N_2\circ f=N_1$, entonces $p\in S_1$ es elíptico $($hiperbólico, llano$)$ si y sólo si $f(p)\in S_2$ es elíptico $($hiperbólico, llano$)$.

        \item[b)] Si $S_1,S_2$ son dos esferas con igual radio $r>0$, $S$ es una superficie regular orientable no contenida en la bola con borde $S_j$ para $j=1,2$, y $S_1\cap S=S_2\cap S=S_1\cap S_2=\{p\}$ para un punto $p\in \mathbb{R}^3$, entonces las curvaturas principales $k_1(p),k_2(p)$ de $S$ en $p$ satisfacen $-1/r\leq k_1(p)\leq k_2(p)\leq 1/r$.

        \item[c)] Si $S_1,S_2$ son superficies regulares orientables, $N_1:S_1\to \mathbb{S}^2$ y $N_2:S_2\to \mathbb{S}^2$ sus aplicaciones de Gauss y $f:S_1\to S_2$ una isometría local, entonces $N_2\circ f=N_1$.

        \item[d)] Existe una parametrización $X:U\to \mathbb{S}^2$ de $\mathbb{S}^2$ tal que $E=G=3$ y $F=0$.
    \end{enumerate}
    \end{ejercicio}

    \newpage
    \setcounter{ejercicio}{0} % Reiniciar contador de ejercicios
    \noindent
    \textbf{Solución.}

    \begin{ejercicio}~
        \begin{enumerate}
            \item[a)] Sea $\alpha : (0,1) \to \mathbb{R}^3$ una curva diferenciable p.p.a. y tal que su traza está
            contenida en la esfera centrada en el origen de radio $r > 0$. Comprueba que la curvatura de $\alpha$ es 
            siempre mayor o igual que $1/r$. \\

            Por definición de $\alpha$ p.p.a, tenemos que $$|\alpha'(t)| = 1 \quad \forall t \in \left]0,1\right[$$
            Por otro lado, como la traza de $\alpha$ está contenida en $S^2(0,r)$, por definición de esta esfera
            \begin{equation}\label{eq:ec1}
                |\alpha(t)|^2 = r^2 \quad \forall t \in \left] 0,1 \right[
            \end{equation}
            Derivamos (\ref{eq:ec1}) en cada miembro respecto de $t$, recordando que 
            $|\alpha(t)|^2 = \prodesc{\alpha(t)}{\alpha(t)}$, y que el producto escalar es simétrico:
            $$2 \prodesc{\alpha(t)}{\alpha'(t)} = \prodesc{\alpha'(t)}{\alpha(t)} + \prodesc{\alpha(t)}{\alpha'(t)} = 
            0 \quad \forall t \in \left] 0,1 \right[$$
            es decir
            \begin{equation}\label{eq:ec2}
                \prodesc{\alpha(t)}{\alpha'(t)} = 0 \quad \forall t \in I
            \end{equation}
            Derivamos otra vez (\ref{eq:ec2}) igual que antes
            $$\overbrace{\prodesc{\alpha'(t)}{\alpha'(t)}}^{(*)} + \prodesc{\alpha(t)}{\alpha''(t)} = 0$$
            y usando que $\alpha$ está p.p.a, $(*) = |\alpha'(t)|^2 = 1$, y
            $$1 + \prodesc{\alpha(t)}{\alpha''(t)} = 0 \iff \prodesc{\alpha(t)}{\alpha''(t)} = -1$$
            Tomando módulos en ambos miembros, y aplicando la desigualdad de Cauchy-Schwartz (CS)
            \begin{equation}\label{eq:ec3}
                1 = |-1| = |\prodesc{\alpha(t)}{\alpha''(t)}| \stackrel{\text{(CS)}}{\leq} |\alpha(t)| |\alpha''(t)|
            \end{equation}
            De (\ref{eq:ec1}) deducimos también, ya que $r > 0$, que $|\alpha(t)| = r$, luego sustituyendo en (\ref{eq:ec3})
            $$1 \leq r |\alpha''(t)| \stackrel{r > 0}{\iff} |\alpha''(t)| \geq \dfrac{1}{r}$$
            Por último, como $\alpha$ está p.p.a. por teoría sabemos que su curvatura en $\R^3$ viene dada por
            $$k_\alpha(t) = |\alpha''(t)|$$
            y es inmediato que
            $$\boxed{k_\alpha(t) = |\alpha''(t)| \geq \dfrac{1}{r}}$$

            \newpage

            \item[b)] Si $\beta : (0,1) \to \mathbb{R}^3$ una curva diferenciable p.p.a. con curvatura no nula en todo 
            instante, entonces $\gamma : (0,1) \to \mathbb{R}^3$, $\gamma(t)=\beta'(t)$, es una curva regular con 
            curvatura mayor o igual que $1$. \\

            Consideramos $I = \left]0,1 \right[$. Por definición de $\gamma$, está claro que $\gamma'(t) = \beta''(t)$. 
            Como $\beta$ está p.p.a., ya hemos visto en el apartado anterior que su curvatura es $k_\beta(t) = |\beta''(t)|$. Por hipótesis
            $$k_\beta(t) \neq 0 \quad \forall t \in I$$
            por lo que 
            $$|\beta''(t)| = k_\beta(t) \neq 0 \Longrightarrow \beta''(t) \neq 0 \Longrightarrow \gamma'(t) \neq 0 \quad \forall t 
            \in I$$
            Por definición entonces $$\boxed{\gamma \text{ es regular.}}$$
            Falta ver que su curvatura es mayor o igual que $1$. Como $\beta$ está p.p.a., 
            $$|\beta'(t)| = 1 \quad \forall t \in I$$ 
            y por definición de $\gamma(t)$
            $$|\gamma(t)| = |\beta'(t)| = 1 \quad \forall t \in I$$
            Es decir, la traza de $\gamma$ está contenida $\Ss^2$. Para poder aplicar el apartado a) con $r = 1$, 
            tenemos que tomar una reparametrización por el arco de $\gamma$, y dicha reparametrización
            si que estará p.p.a. (con $\gamma$ no lo podemos demostrar por falta de hipótesis). Fijado $a \in 
            I$, consideramos la aplicación longitud orientada de $\gamma$ respecto al instante $a$
            vista en teoría
            \Func{S_a}{I}{\R}{t}{S_a(t) := \int_{a}^{t} |\gamma'(u)| du}
            Como $\gamma$ es regular, sabemos que $S_a : I \to S_a(I) = J$ es un difeomorfismo sobreyectivo, por
            lo que considerando $h = (S_a)^{-1} : J \to I$ obtenemos la reparametrización $\tilde{\gamma} = \gamma 
            \circ h : J \to \R^3$. Por definición, $S_a(h(s)) = s$, y derivando $$1 = h'(s)S_a'(h(s))$$
            Como $S'_a(t) = |\gamma'(t)|$ por el Teorema Fundamental del Cálculo, $$1 = h'(s)|\gamma'(h(s))| \Longrightarrow
            h'(s) = \dfrac{1}{|\gamma'(h(s))|}$$
            Derivando $\tilde{\gamma}$ con la regla de la cadena
            $$\tilde{\gamma}'(s) = h'(s) \gamma'(h(s)) = \dfrac{\gamma'(h(s))}{|\gamma'(h(s))|} \Longrightarrow
            |\tilde{\gamma}'(s)| = 1 \quad \forall s \in J$$
            por lo que $\tilde{\gamma}$ está p.p.a.. Lo importante es que, al ser $\tilde{\gamma}$ una 
            reparametrización de $\gamma$, se tiene que ambas trazas son iguales, $\tr(\tilde{\gamma}) = \tr(\gamma)$,
            y como ya sabíamos que $\tr(\gamma) \subset \Ss^2$, entonces
            $$\tr(\tilde{\gamma}) = \tr(\gamma) \subset \Ss^2$$
            y ahora sí puede aplicarse el apartado a) con $r=1$ a $\tilde{\gamma}$ para obtener que 
            $$k_{\tilde{\gamma}}(s) \geq 1 \quad \forall s \in J$$
            En particular, para cada $t \in I$ como $S_a(t) \in J$, se tiene que
            $$k_{\tilde{\gamma}}(S_a(t)) \geq 1$$
            y por definición de la curvatura de $\gamma$ según su reparametrización, 
            $k_\gamma(t) = k_{\tilde{\gamma}}(S_a(t))$, luego
            $$\boxed{k_\gamma(t) \geq 1 \quad \forall t \in I}$$

        \end{enumerate}
    \end{ejercicio}

    \newpage

    \begin{ejercicio}
        Sea $S=\{X(u,v)\mid u\in (0,+\infty), v\in (-\frac{\pi}{2},\frac{\pi}{2})\}$ donde
        $$
        X(u,v)=(u\cos(v),u\sen(v),\ln(\cos(v))+u).
        $$

        Se pide:
        \begin{enumerate}
            \item[a)] Demuestra que $X$ es una parametrización global de $S$ y que $S$ es una superficie orientable. \\
            
            Vamos por orden. Primero obtenemos una expresión en función de $(x,y,z) \in \R^3$ de $S$. Expresando $(x,y,z) = X(u,v)$, 
            entonces
            $$x = u \cos v, \quad y = u \sen v$$
            Como $u > 0$ y $\cos v > 0$ (ya que $v \in (-\frac{\pi}{2},\frac{\pi}{2})$), $x > 0$. Además,
            identificando el cambio de coordenadas cartesianas a polares, tenemos que $u = \sqrt{x^2 + y^2}$. También
            $$\cos v = \dfrac{x}{u} = \dfrac{x}{\sqrt{x^2 + y^2}}$$ por lo que
            $$z = \ln(\cos v) + u = \ln \left( \dfrac{x}{\sqrt{x^2 + y^2}} \right) + \sqrt{x^2 + y^2}$$
            Entonces, los puntos $(x,y,z) \in S$ verifican
            $$z - \sqrt{x^2 + y^2} - \ln \left( \dfrac{x}{\sqrt{x^2 + y^2}} \right) = 0$$
            Vemos ahora que $S$ es una superficie (regular), por ser imagen inversa (no vacía) del valor regular $0$, 
            por la aplicación diferenciable \Func{f}{O}{\R}{(x,y,z)}{z - \sqrt{x^2 + y^2} - \ln \left( \dfrac{x}{\sqrt{x^2 + y^2}} \right)} 
            donde $O = \{(x,y,z) \in \R^3 : x > 0\}$ (abierto por ser preimagen
            de un abierto ($\Rmas$) por la proyección $(x,y,z) \mapsto x$ (continua)). Entonces $\emptyset \neq S = f^{-1}(\{0\})$
            (por ejemplo, $(1, 0, 1) = X(1,0) \in S$). Comprobamos que $0$ es un valor regular
            por definición, pues basta observar que $$\dfrac{\partial f}{\partial z}(x,y,z) = 1$$
            para $(x,y,z) \in S$, luego 
            $$\nabla f(x,y,z) \neq 0 \quad \forall (x,y,z) \in S$$
            Veamos ahora que $S$ es orientable, lo cual es sencillo teniendo en cuenta que la superficie viene dada como imagen inversa de un valor 
            regular. Sabemos por teoría que la forma natural de definir la aplicación de Gauss que determina la orientación de $S$ 
            (denotada por $N_0$ pensando en lo que sigue) es:
            \Func{N_0}{S}{\R^3}{(x,y,z)}{N_0(x,y,z) = \dfrac{\nabla f}{|\nabla f|}(x,y,z)}
            El denominador no se anula nunca por ser $S$ superficie regular, luego $N_0$ está definida en todo $S$, es diferenciable (tipo a)) y, 
            por definición, $S$ es orientable. \\

            Falta por ver que $X$ es una parametrización global de $S$, para lo cual, primero vemos que es local. Denotamos 
            $$U = \Rmas \times \left] -\dfrac{\pi}{2}, \dfrac{\pi}{2} \right[ $$
            Probamos primero que $X: U \to \R^3$ dada en el enunciado es una parametrización local de $S$, según la definición vista en teoría.
            Lo primero es que el dominio de la parametrización $U$ es abierto en $\R^2$ por 
            ser producto de abiertos de $\R$. Las componentes de $X$ son $u \cos v$, $u \sen v$ y $\ln(\cos v) + u$, todas ellas 
            funciones diferenciables de $(u, v)$ en $U$, ya que si $$v \in \left] -\dfrac{\pi}{2}, \dfrac{\pi}{2} \right[$$ 
            entonces $\cos v > 0$, por lo que $\ln(\cos v)$ está bien definido. Entonces, $X$ es diferenciable (en el sentido del análisis). \\

            Ahora, obtenemos $X(U)$, con el fin a su vez de determinar el entorno parametrizado $V$ correspondiente por $X$. Por 
            definición, $S = \{X(u,v) : (u,v) \in U\}$, luego $X(U) = S$, y tomamos $V = S$. $V$ es un abierto de $S$
            (es el total), luego es un entorno abierto de cualquier punto de $S$. 
            Falta probar que $X : U \to V$ es un homeomorfismo y que $(dX)_q : \R^2 \to \R^3$
            es inyectiva para todo $q \in U$. \\

            Para lo primero, ya sabemos que $X : U \to V$ es sobreyectiva por definición de $V$. Debemos ver que $X$ es inyectiva, y que 
            $X^{-1} : V \to U$ es continua. La inyectividad se deduce fácilmente, pues, si $X(u, v) = X(\rho, \theta)$, 
            con $(u, v), (\rho, \theta) \in U$, entonces, coordenada
            a coordenada
            $$(u \cos v, u \sen v, \ln(\cos v) + u) = X(u,v) = X(\rho, \theta) = (\rho \cos \theta, \rho \sen \theta, \ln(\cos \theta) + \rho) 
            \Longrightarrow$$
            $$\begin{cases}
                u \cos v = \rho \cos \theta \\
                u \sen v = \rho \sen \theta \\
                \ln(\cos v) + u = \ln(\cos \theta) + \rho
            \end{cases}$$
            
            Elevando al cuadrado las dos primeras igualdades, y sumándolas, obtenemos
            $$u^2(\cos^2 v + \sen^2 v) = (u \cos v)^2 + (u \sen v)^2 = (\rho \cos \theta)^2 + (\rho \sen \theta)^2 = \rho^2(\cos^2 \theta + \sen^2 \theta)$$
            de donde
            $$u^2 = \rho^2 \stackrel{u,\rho>0}{\Longrightarrow} u = \rho$$
            Entonces, en las dos primeras ecuaciones pueden
            simplificarse las amplitudes de los senos y cosenos, quedando
            $$\begin{cases}
                \cos v = \cos \theta \\
                \sen v = \sen \theta
            \end{cases}$$
            de donde $v - \theta \in 2 \pi \Z$, pero $v, \theta \in \left]- \frac{\pi}{2} , \frac{\pi}{2} \right[$, por lo que 
            necesariamente $v = \theta$, y a su vez, $(u, v) = (\rho, \theta)$, luego $X$ es inyectiva. \\

            La aplicación inversa estará dada por 
            \Func{X^{-1}}{V}{U}{(x,y,z)}{X^{-1}(x,y,z) = (u,v)}
            La primera componente de la aplicación inversa ya sabemos que es $u = \sqrt{x^2 + y^2}$, con $u > 0$ ya que $(x,y,z) \in S$,
            luego $x > 0$. Para la segunda, correspondiente con $v$, usamos nuevamente que $(x,y,z)$ es un punto de $V = S = X(U)$. 
            Es decir, existe $(u,v) \in U$ tal que $(x,y,z) = X(u,v)$. Como $u > 0$, y $x > 0$, entonces, como $x = u \cos v$,
            debe ser $\cos v > 0$. Así, como $y = u \sen v$, entonces
            $$\dfrac{y}{x} = \dfrac{u \sen v}{u \cos v} \stackrel{u>0}{=} \tan v$$
            Como $v \in \left]- \frac{\pi}{2} , \frac{\pi}{2} \right[$, deducimos que
            $$v = \arctan \left( \dfrac{y}{x} \right)$$
            Al ser composición de funciones continuas cada componente (la segunda es continua ya que $x > 0$), $X^{-1}$ es continua. 
            Por lo tanto, $X : U \to V$ es un homeomorfismo. \\

            Falta por ver que la diferencial es inyectiva, lo cual sabemos por teoría que equivale a ver que las parciales $X_u$ y 
            $X_v$ son linealmente independientes. En este caso, dado que $X(u,v)=(u\cos(v),u\sen(v),\ln(\cos(v))+u)$, calculamos
            $$X_u = (\cos v, \sen v, 1)$$
            $$X_v = (-u \sen v, u \cos v, -\tan v)$$
            y como 
            $$\begin{vmatrix}
                i & j & k \\
                \cos v & \sen v & 1 \\
                - u \sen v & u \cos v & - \tan v
            \end{vmatrix} = \left(-u \cos v - \dfrac{\sen^2 v}{\cos v}, (1-u) \sen v, u \right)$$

            y $u > 0$, se deduce que $X_u(u, v) \land X_v(u, v) \neq 0 \quad \forall (u, v) \in U$. \\
            
            En definitiva, hemos demostrado que $X : U \to \R^3$ es una parametrización local de $S$, donde el entorno parametrizado es 
            $V = S$. Como además $X(U) = S$, esta parametrización local cubre toda la superficie, por lo que 
            $$\boxed{X \text{ es una parametrización global de } S.}$$

            \item[b)] Calcula la primera y segunda formas fundamentales de $S$. Prueba que el punto $p=(1,0,1)\in S$ y determina el plano 
            tangente de la superficie en dicho punto. \\

            Trabajamos en la base del plano tangente $B = \{X_u, X_v\}$. La primera forma fundamental sabemos que se define,
            para $p = X(u,v) \in S$, como \Func{g_p}{T_p C \times T_p C}{\R}{(v,w)}{g_p(v,w) := \prodesc{v}{w}}

            Ahora, obtenemos igual que en los ejemplos vistos en teoría los coeficientes de la primera forma fundamental
            $$E = \prodesc{X_u}{X_u} = |X_u|^2 = \cos^2 v + \sen^2 v + 1 = 2$$
            $$F = \prodesc{X_u}{X_v} = - u \cos v \sen v + u \cos v \sen v - \tan v = - \tan v$$
            $$G = \prodesc{X_v}{X_v} = |X_v|^2 = u^2 \sen^2 v + u^2 \cos^2 v + \tan^2 v = u^2 + \tan^2 v$$

            Por lo tanto, la matriz de la primera forma fundamental $g_p$ en la base $B$ es 
            $$\boxed{M_B(g_p) = \begin{pmatrix}
                E & F \\
                F & G
            \end{pmatrix} = \begin{pmatrix}
                2 & -\tan v \\
                -\tan v & u^2 + \tan^2 v
            \end{pmatrix}}$$

            Nuevamente, para $p = X(u,v) \in S$, recordamos la definición de la segunda forma fundamental
            \Func{\sigma_p}{T_p C \times T_p C}{\R}{(v,w)}{\sigma_p(v,w) := \prodesc{A_p(v)}{w}}

            donde $A_p$ es el endomorfismo de Weingarten. Para determinarla, usaremos las fórmulas de los ejemplos de teoría, para 
            lo cual, dado que $X$ es una parametrización global de $S$, fijamos la orientación dada por
            $$N^X = N \circ X = \dfrac{X_u \land X_v}{|X_u \land X_v|}$$
            Obtenemos el denominador (ya que el numerador lo hemos calculado antes) usando una fórmula vista en teoría 
            $$|X_u \land X_v|^2 = EG - F^2 = 2(u^2 + \tan^2 v) - (- \tan v)^2 = 2u^2 + \tan^2 v \Longrightarrow$$
            $$|X_u \land X_v| = \sqrt{2u^2 + \tan^2v}$$

            Por tanto, $$N^X = \dfrac{X_u \land X_v}{|X_u \land X_v|} = 
            \dfrac{1}{\sqrt{2u^2 + \tan^2v}} \left(-u \cos v - \dfrac{\sen^2 v}{\cos v}, (1-u) \sen v, u \right)$$
            Obtenemos ahora las derivadas segundas de $X$, recordando primero que
            $$X_u = (\cos v, \sen v, 1)$$
            $$X_v = (-u \sen v, u \cos v, -\tan v)$$
            Entonces
            $$X_{uu} = (0,0,0)$$
            $$X_{uv} = X_{vu} = (-\sen v, \cos v, 0)$$
            $$X_{vv} = (-u \cos v, - u \sen v, -\sec^2 v)$$

            Calculamos ahora los coeficientes de la segunda forma fundamental como se ha visto en los ejemplos de teoría
            $$e = \prodesc{X_{uu}}{N^X} = \prodesc{(0,0,0)}{N^X} = 0$$
            $$f = \prodesc{X_{uv}}{N^X} = \dfrac{\det(X_u, X_v, X_{uv})}{|X_u \land X_v|} = 
            \dfrac{1}{\sqrt{2u^2 + \tan^2 v}} \begin{vmatrix}
                \cos v & \sen v & 1 \\
                - u \sen v & u \cos v & - \tan v \\
                - \sen v & \cos v & 0
            \end{vmatrix} = $$
            $$\dfrac{1}{\sqrt{2u^2 + \tan^2 v}} (\cancelto{0}{-u \sen v \cos v + u \cos v \sen v} + \tan v (\cos^2 v + \cos^2 v) + 0) = 
            \dfrac{\tan v}{\sqrt{2u^2 + \tan^2 v}}$$
            $$g = \prodesc{X_{vv}}{N^X} = \dfrac{\det(X_u, X_v, X_{vv})}{|X_u \land X_v|} = 
            \dfrac{1}{\sqrt{2u^2 + \tan^2 v}} \begin{vmatrix}
                \cos v & \sen v & 1 \\
                - u \sen v & u \cos v & - \tan v \\
                - u \cos v & - u \sen v & -\sec^2 v
            \end{vmatrix} = $$
            $$\dfrac{1}{\sqrt{2u^2 + \tan^2 v}} (u^2(\sen^2 v + \cos^2 v) - \tan v\cancelto{0}{(-u \cos v \sen v + u \cos v \sen v)} + 
            u \sec^2) = $$
            $$\dfrac{u^2 - u \sec^2 v}{\sqrt{2u^2 + \tan^2 v}}$$

            Ahora, la matriz de la segunda forma fundamental $\sigma_p$ en la base $B$ es
            $$\boxed{M_B(\sigma_p) = \begin{pmatrix}
                e & f \\
                f & g
            \end{pmatrix} = \dfrac{1}{\sqrt{2u^2 - \tan^2v}} \begin{pmatrix}
                0 & \tan v \\
                \tan v & u^2 - u \sec^2v
            \end{pmatrix}}$$

            Quedan por ver las dos últimas cosas que dice el apartado. Buscamos $(u,v) \in U$ tal que $(1,0,1) = p = X(u,v)$. 
            Recordemos que $$X(u,v)=(u\cos(v),u\sen(v),\ln(\cos(v))+u).$$ Entonces, para la primera y segunda componente por ejemplo,
            imponemos que 
            $$\begin{cases}
                u \cos v = 1 \\
                u \sen v = 0
            \end{cases} \stackrel{u>0}{\Longrightarrow} \sen v = 0 \Longrightarrow v = 0$$

            ya que $v \in \left] - \dfrac{\pi}{2}, \dfrac{\pi}{2} \right[$. Sustituyendo en la primera ecuación, $u = u \cos 0 = 1$. \\

            En efecto, $X(u,v) = X(1,0) = (1, 0, 1) = p$. Determinamos ahora el plano tangente a $S$ en $p$. Como $p = X(1,0)$, 
            una base de $T_pS$ viene dada por $\{X_u(1,0), X_v(1,0)\}$. Evaluamos
            $$X_u = (\cos v, \sen v, 1)$$
            $$X_v = (-u \sen v, u \cos v, -\tan v)$$
            en $(1,0)$ y determinamos la base
            $$X_u(1,0) = (1,0,1)$$
            $$X_v(1,0) = (0,1,0)$$
            Así, $T_pS = \langle (1,0,1), (0,1,0) \rangle$. Para dar la ecuación cartesiana del plano, obtenemos el vector normal
            $$\begin{vmatrix}
                i & j & k \\
                1 & 0 & 1 \\
                0 & 1 & 0
            \end{vmatrix} = (-1,0,1)$$

            La ecuación del plano $T_pS$ es $Ax + By + Cz + D = 0$, con $(A,B,C) = (-1,0,1)$ el vector normal que hemos calculado. Sustituyendo
            y usando que $p \in T_p S$, obtenemos $D$
            $$-x + z + D = 0 \stackrel{p = (1,0,1)}{\Longrightarrow} -1 + 1 + D = 0 \Longrightarrow D = 0$$
            Y entonces
            $$\boxed{T_p S = \{(x,y,z) \in \R^3 : -x + z = 0 \} = \{(x,y,z) \in \R^3 : x = z \}}$$

            \item[c)] Calcula la curvatura de Gauss y la curvatura media de $S$ y clasifica los puntos de $S$ en elípticos, hiperbólicos, 
            llanos $($parabólicos y planos$)$. \\

            Para obtener $K$ y $H$ curvatura de Gauss y curvatura media de $S$, respectivamente, usamos las fórmulas vistas en teoría,
            para lo cual consideramos las expresiones locales de las curvaturas en la parametrización $X$: $K_X = K \circ X$ y 
            $H_X = H \circ X$, y entonces
            $$K_X = \dfrac{eg - f^2}{EG - F^2}$$
            $$H_X = \dfrac{eG - 2fF + gE}{2(EG-F^2)}$$
            Sustituyendo los valores del apartado anterior previamente calculados
            $$E = 2, \quad F = - \tan v, \quad G = u^2 + \tan^2 v, \quad e = 0, \quad f = \dfrac{\tan v}{\sqrt{2u^2 + \tan^2v}},
            \quad g = \dfrac{u^2 - u\sec^2 v}{\sqrt{2u^2 + \tan^2v}}$$
            nos queda
            $$K_X(u,v) = \dfrac{-f^2}{EG-F^2} = \dfrac{-\dfrac{\tan^2v}{2u^2 + \tan^2v}}{2u^2 + \tan^2v} = -\dfrac{\tan^2v}{(2u^2 + \tan^2v)^2}$$
            y para la curvatura media, el numerador es
            $$-2fF + gE = \dfrac{2\tan^2 v}{\sqrt{2u^2 + \tan^2v}} + \dfrac{2(u^2 - u \sec^2v)}{\sqrt{2u^2 + \tan^2v}} = 
            \dfrac{2\tan^2 v + 2u^2 - 2u \sec^2 v}{\sqrt{2u^2 + \tan^2v}}$$
            de donde
            $$H_X(u,v) = \dfrac{2 \tan^2v + 2u^2 - 2u \sec^2 v}{2(2u^2 + \tan^2v)\sqrt{2u^2 + \tan^2v}} = 
            \dfrac{u^2 - u \sec^2v + \tan^2 v}{(2u^2 + \tan^2 v)^{3/2}}$$
            Ahora, como $X$ es una parametrización global, para todo punto $p = X(u,v) \in S$, se tiene
            $$K(p) = K_X(u,v) \quad H(p) = H_X(u,v)$$
            Es decir,
            $$\boxed{K(X(u,v)) = -\dfrac{\tan^2v}{(2u^2 + \tan^2v)^2} \quad \forall (u,v) \in U}$$
            $$\boxed{H(X(u,v)) = \dfrac{u^2 - u \sec^2v + \tan^2 v}{(2u^2 + \tan^2 v)^{3/2}} \quad \forall (u,v) \in U}$$

            Respecto a la clasificación de los puntos, no podemos decir nada porque no se ha definido qué significa cada tipo punto en teoría. 

            \item[d)] Prueba que $S$ es una superficie reglada. Razonar si $S$ es localmente isométrica o no a un plano. \\
            
            No se responde dado que ni siquiera se ha definido en teoría qué es una superficie reglada, ni una isometría y menos aún una 
            isometría local.

            \item[e)] Sea $\pi:S\to \{z=0\}$ la aplicación dada por $\pi(x,y,z)=(x,y,0)$. Prueba que $\pi$ es un difeomorfismo 
            sobre su imagen. \\

            Tenemos que ver que $\pi : S \to \pi(S)$ es biyectiva, $\pi$ es diferenciable, y su inversa $\pi^{-1} : \pi(S) \to S$
            también es diferenciable. \\

            Primero, obtenemos $\pi(S)$. Como $S = X(U)$, tenemos que $$\pi(X(u,v)) = (u \cos v, u \sen v, 0)$$
            Ahora, si $(x,y,0) \in \pi(S)$, entonces existen $(u,v) \in U$ tales que
            $$(x,y,0) = (u \cos v, u \sen v, 0)$$
            como $\cos v > 0$ porque $v \in \left] - \frac{\pi}{2}, \frac{\pi}{2} \right[$, entonces $x = u \cos v > 0$.
            Así $$\pi(S) \subseteq \{(x,y,0) \in \R^3 : x > 0\}$$
            Recíprocamente, si $(x,y,0) \in R^3$ con $x >0$, tomando los valores que ya comprobamos en la demostración 
            de que $X^{-1}$ era continua, $$u = \sqrt{x^2 + y^2}, \quad v = \arctan \left( \dfrac{y}{x} \right)$$
            Entonces $u>0$, y $$v \in \left] - \frac{\pi}{2}, \frac{\pi}{2} \right[$$
            Además, $u \cos v = x$ y $u \sen v = y$, luego
            $$(x,y,0) = \pi(X(u,v)) \in \pi(S)$$
            por lo que
            $$\pi(S) = \{(x,y,0) \in \R^3 : x > 0\}$$
            Ahora, como para cada $(x,y,0) \in \pi(S)$ los valores que hemos tomado de $u$ y $v$ son únicos
            (dado que $x > 0$, y $v \in \left]-\frac{\pi}{2}, \frac{\pi}{2} \right[$, donde $\tan v$ es inyectiva),
            $\pi : S \to \pi(S)$ es biyectiva. Su inversa es
            \Func{\pi^{-1}}{\pi(S)}{S}{(x,y,0)}{\pi^{-1}(x,y,0) = X\left(\sqrt{x^2 + y}, \arctan \left( \dfrac{y}{x} \right)\right)}
            es decir
            $$\pi^{-1}(x,y,0) = \left(x, y, \sqrt{x^2 + y^2} + \ln \left( \dfrac{x}{\sqrt{x^2 + y^2}} \right)\right)$$
            La aplicación $\pi$ es diferenciable, como restricción a $S$ de la aplicación diferenciable $(x,y,z) \mapsto (x,y,0)$. 
            Además, $\pi^{-1}$ es diferenciable por serlo sus componentes en el abierto $x > 0$ (ya se comprobó esto en el principio
            del apartado 2a)). Por lo tanto, 
            $$\boxed{\pi : S \to \pi(S) \text{ es un difeomorfismo.}}$$
        \end{enumerate}
    \end{ejercicio}

    \newpage

    \begin{ejercicio}
        Razonar si son verdaderas o falsas las siguientes afirmaciones:

        \begin{enumerate}
            \item[a)] Si $S_1,S_2$ son superficies regulares orientables con aplicaciones de Gauss $N_1:S_1\to \mathbb{S}^2$ y $N_2:S_2\to 
            \mathbb{S}^2$ y $f:S_1\to S_2$ es un difeomorfismo local tal que $N_2\circ f=N_1$, entonces $p\in S_1$ es elíptico $($hiperbólico, 
            llano$)$ si y sólo si $f(p)\in S_2$ es elíptico $($hiperbólico, llano$)$. \\

            No se responde al no haberse definido qué son los puntos elípticos, hiperbólicos, ni llanos.

            \item[b)] Si $S_1,S_2$ son dos esferas con igual radio $r>0$, $S$ es una superficie regular orientable no contenida en la bola 
            con borde $S_j$ para $j=1,2$, y $S_1\cap S=S_2\cap S=S_1\cap S_2=\{p\}$ para un punto $p\in \mathbb{R}^3$, entonces las 
            curvaturas principales $k_1(p),k_2(p)$ de $S$ en $p$ satisfacen $-1/r\leq k_1(p)\leq k_2(p)\leq 1/r$. \\

            $\boxed{\text{Verdadera}}$. Por hipótesis, $S$ es una superficie orientable, luego elegimos una orientación de $S$,
            de tal manera que $S$ queda orientada, y denotamos por $N(p)$ a la normal unitaria en $p$. Como las dos esferas
            tienen el mismo radio $r>0$, y son tangentes en $p$, sus centros son $$c_{+} = p + rN(p), \quad c_{-} = p - rN(p)$$
            Tomamos $v \in T_pS$ un vector unitario, y por definición de vector en el plano tangente, sea $\alpha$ una curva p.p.a.
            en $S$ tal que $$\alpha(0) = p \quad \alpha'(0) = v$$
            Consideramos primero la esfera de centro $c_{+}$. Definimos
            $$\varphi_{+}(t) = |\alpha(t) - c_{+}|^2 - r^2$$
            Como $S$ no está contenida en la bola con borde $S_j$, $j = 1,2$, y solo toca a la esfera en $p$, entonces
            $t = 0$ es un mínimo local de $\varphi_{+}$. Por lo tanto, $\varphi''_{+}(0) \geq 0$. Obtenemos $\varphi_+'(t)$ y $\varphi_+''(t)$
            $$\varphi_{+}'(t) = 2 \prodesc{\alpha(t) - c_{+}}{\alpha'(t)}$$
            y 
            $$\varphi_{+}''(t) = 2(\prodesc{\alpha'(t)}{\alpha'(t)} + \prodesc{\alpha(t) - c_{+}}{\alpha''(t)}) = 
            2|\alpha'(t)|^2 + 2 \prodesc{\alpha(t) - c_{+}}{\alpha''(t)}$$
            Como $\alpha$ está p.p.a, $|\alpha'(0)| = 1$. Además
            $$\alpha(0) - c_{+} = p - (p + rN(p)) = - r N(p)$$
            Por lo que $$0 \leq \varphi_{+}''(0) = 2 - 2r\prodesc{\alpha''(0)}{N(p)} \quad(*)$$
            En este punto, aparece la segunda forma fundamental, pero está implícita en el último producto escalar. 
            Como $N(\alpha(t))$ es normal a $S$ y $\alpha'(t)$ es tangente a $S$, se tiene que
            $$\prodesc{N(\alpha(t))}{\alpha'(t)} = 0$$
            Derivando y evaluando en $t=0$
            $$\prodesc{(dN)_p(\alpha'(0))}{\alpha'(0)} + \prodesc{N(p)}{\alpha''(0)} = 0$$
            Como $\alpha'(0) = v$, entonces
            $$\prodesc{(dN)_p(v)}{v} + \prodesc{N(p)}{\alpha''(0)} = 0 \Longrightarrow -\prodesc{(dN)_p(v)}{v} = \prodesc{\alpha''(0)}{N(p)}$$
            Por definición de la segunda forma fundamental $$\sigma_p(v,v) = \prodesc{\alpha''(0)}{N(p)}$$
            y se tiene (despejando en $(*)$) que $$\sigma_p(v,v) \leq \dfrac{1}{r}$$
            De manera análoga se llega a que $$\sigma_p(v,v) \geq -\dfrac{1}{r}$$
            usando la otra esfera de centro $c_{-} = p - rN(p)$. Así, hemos probado que para toda dirección unitaria 
            $v \in T_p S$, 
            $$-\dfrac{1}{r} \leq \sigma_p(v,v) \leq \dfrac{1}{r}$$
            Por un ejercicio hecho en teoría, sabemos que 
            \begin{equation*}
                k_1(p) = \min_{\substack{v\in T_pS\\|v| =1}} \sigma_p(v,v), \qquad k_2(p) = \max_{\substack{v\in T_pS\\|v| =1}}\sigma_p(v,v)
            \end{equation*}
            por lo que en particular, para el convenio $k_1(p) \leq k_2(p)$:
            $$\boxed{-\dfrac{1}{r} \leq k_1(p) \leq k_2(p) \leq \dfrac{1}{r}}$$ 
            

            \item[c)] Si $S_1,S_2$ son superficies regulares orientables, $N_1:S_1\to \mathbb{S}^2$ y $N_2:S_2\to \mathbb{S}^2$ sus aplicaciones 
            de Gauss y $f:S_1\to S_2$ una isometría local, entonces $N_2\circ f=N_1$. \\

            No se responde por el mismo motivo que en 2d).

            \item[d)] Existe una parametrización $X:U\to \mathbb{S}^2$ de $\mathbb{S}^2$ tal que $E=G=3$ y $F=0$. \\
            
            $\boxed{\text{Falsa}}$. Por reducción al absurdo, supongamos que existe una parametrización en las condiciones del enunciado.
            Como $X(U) \subset \Ss^2$, entonces $$|X(u,v)|^2 = \prodesc{X(u,v)}{X(u,v)} = 1 \quad \forall (u,v) \in U$$ Derivando 
            respecto de $u$ y de $v$ la última igualdad, obtenemos
            $$2 \prodesc{X_u}{X} = \prodesc{X_u}{X} + \prodesc{X}{X_u} = \dfrac{\partial}{\partial u} \prodesc{X}{X} = 0 \Longrightarrow
            $$
            \begin{equation}\label{eq:ec4}
                \prodesc{X_u}{X} = 0
            \end{equation}
            $$2 \prodesc{X_v}{X} = \prodesc{X_v}{X} + \prodesc{X}{X_v} = \dfrac{\partial}{\partial v} \prodesc{X}{X} = 0 \Longrightarrow 
            \prodesc{X_v}{X} = 0$$
            Además, sabemos que los coeficientes de la primera forma fundamental son
            $$\prodesc{X_u}{X_u} = E = 3, \quad \prodesc{X_u}{X_v} = F = 0, \quad \prodesc{X_v}{X_v} = G = 0$$
            Derivando (\ref{eq:ec4}) respecto de $u$
            $$\prodesc{X_{uu}}{X} + \prodesc{X_u}{X_u} = 0 \Longrightarrow \prodesc{X_{uu}}{X} = -3$$
            Como $E$ y $F$ son constantes, derivando sus expresiones se tiene que
            $$\prodesc{X_{uu}}{X_u} = 0 \quad \prodesc{X_u}{X_{uv}} = 0$$
            Por tanto, como $B = \{X, X_u, X_v\}$ es una base de $\R^3$, está claro que
            \begin{equation}\label{eq:ec5}
                X_{uu} = -3X
            \end{equation}
            Por otro lado, si derivamos (\ref{eq:ec4}) respecto de $v$
            $$\prodesc{X_{uv}}{X} + \prodesc{X_u}{X_v} = 0$$
            y como $F=0$,
            $$\prodesc{X_{uv}}{X} = 0$$
            Nuevamente, como $E$ y $G$ son constantes, derivando sus expresiones respecto de $v$
            $$\prodesc{X_{uv}}{X_u} = 0, \quad \prodesc{X_{uv}}{X_v} = 0$$
            Y de nuevo en coordenadas de la base $B$ tenemos que $X_{uv} = 0$. Entonces $X_{uuv} = 0$, pero derivando (\ref{eq:ec5})
            respecto de $v$ nos queda que $0 = X_{uuv} = -3 X_v$, luego necesariamente $X_v = 0$, pero esto contradice que
            $$|X_v|^2 = \prodesc{X_v}{X_v} = G = 3$$
            Por lo tanto, tal parametrización no puede existir.
        \end{enumerate}
    \end{ejercicio}

\end{document}
